\subsection*{Overall Architecture}
The proposed QDCR-Net contains five stages: dual-image input, dual-branch feature
extraction, bidirectional cross-residual interaction, quality-aware fusion, and
small-object-oriented detection. The raw branch receives the original RGB image, while the
enhanced branch receives a lightweight enhancement result generated by a unified preprocessing
pipeline.

\subsection*{Dual-Branch Feature Extraction}
Let $I_r$ and $I_e$ denote the raw and enhanced inputs. Their shallow encoders are separated
to preserve branch-specific low-level statistics:
\begin{align}
F_r &= E_r(I_r), \\
F_e &= E_e(I_e).
\end{align}
The separation is particularly important in early stages because raw and enhanced images
often differ most strongly in texture, edge strength, and local contrast.

\subsection*{Cross-Residual Interaction}
To allow explicit bidirectional compensation, we introduce a cross-residual interaction
module at selected backbone stages:
\begin{align}
\tilde{F}_r &= F_r + \alpha \cdot \phi(F_e), \\
\tilde{F}_e &= F_e + \beta \cdot \psi(F_r),
\end{align}
where $\phi(\cdot)$ and $\psi(\cdot)$ are $1 \times 1$ alignment convolutions, and
$\alpha$, $\beta$ are learnable gates. The module enables the raw branch to provide stable
semantic priors and the enhanced branch to compensate structural detail responses.

\subsection*{Quality-Aware Dynamic Fusion}
We further introduce a lightweight quality estimator that produces a quality descriptor
$q$ from the input image. The descriptor is mapped to an adaptive fusion weight:
\begin{align}
w = \sigma(W_q q),
\end{align}
and the fused representation is written as:
\begin{align}
F_{fuse} = w \odot \tilde{F}_r + (1 - w) \odot \tilde{F}_e.
\end{align}
This design allows the detector to rely more on raw semantics or enhanced details according
to image quality and degradation characteristics.

\subsection*{Small-Object-Oriented Neck}
To strengthen the localization of weak and tiny underwater targets, we append a shallow
detail path that injects higher-resolution features into the neck. The path is designed to
improve small-object sensitivity while keeping parameter growth limited.
